计算 \(\int_0^1\int_0^1 (x^2+y^2)\,dx\,dy\)。按 Fubini 定理，因被积函数非负，可先积 \(x\)：\(\int_0^1(x^2+y^2)\,dx=[\frac{x^3}{3}+y^2x]_0^1=\frac13+y^2\)，再积 \(y\)：\(\int_0^1(\frac13+y^2)\,dy=\frac13+\frac13=\frac23\)。Fubini 的实质是：非负或绝对可积时，二重积分等于两次单积分的迭代。

前面各单元处理的是单个空间上的测度与积分。一旦两个随机量同时出现——两次测量、输入与标签、状态与观测——问题立刻变成：联合规律定义在哪个空间上，``独立''到底是什么意思，以及二重积分什么时候能拆成两次单积分。

本单元回答三个问题：

\begin{enumerate}
\tightlist
\item
  为什么独立性需要乘积测度：联合事件生活在乘积空间上，``独立''被严格写成联合测度等于边缘测度的乘积。
\item
  Fubini 与 Tonelli 把二重积分变成两次积分：交换积分顺序的两个定理各自要求什么，正负抵消时交换为什么会改答案。
\item
  独立随机变量与联合密度：独立性怎样推出联合密度相乘、和的分布是卷积，以及``不相关但不独立''在测度上意味着什么。
\end{enumerate}

多维对象的规则由一维对象组装而成，但交换积分次序需要单独验证。独立性与可交换性都是因式分解，分解失败时，直觉的``先积一个变量再积另一个''就可能给出错误答案。

\subsection{一个具体算例}

Tonelli 适用（非负）：\(\int_0^1\int_0^1 \frac{1}{(1+xy)^2}\,dx\,dy\)，先积 \(x\) 得 \(\int_0^1\frac{1}{1+y}\,dy=\ln 2\)。

Fubini 需绝对可积：\(\int_0^\infty\int_0^\infty e^{-(x+y)}\,dx\,dy=1\)，交换次序结果相同。但 \(\int_1^\infty\int_0^\infty \frac{x^2-y^2}{(x^2+y^2)^2}\,dx\,dy\) 不绝对可积，交换次序得到 \(\pi/2\) vs \(-\pi/2\)——答案变了。

\subsection{怎么读这一章}

先用乘积测度说明怎样从两个空间构造联合空间，再看 Tonelli 与 Fubini 在什么条件下允许改变积分次序，最后回到随机变量的联合分布与独立性。尤其要分清非负函数与绝对可积函数对应的不同定理。

\subsection{本章篇目}

以下篇目按概念依赖排列；若从具体问题进入，也可以先打开最接近当前困惑的一篇，再沿篇内链接回补。

\begin{itemize}
\tightlist
\item \hyperref[sec:12-Real-Analysis-Support/06-Product-Measures-and-Fubini/01]{``为什么独立性需要乘积测度''}；
\item \hyperref[sec:12-Real-Analysis-Support/06-Product-Measures-and-Fubini/02]{``Fubini 与 Tonelli：二重积分何时能拆成两次''}；
\item \hyperref[sec:12-Real-Analysis-Support/06-Product-Measures-and-Fubini/03]{``独立随机变量与联合密度''}。
\end{itemize}
